\documentclass[exercice,a4]{thedocument}%
\usepackage{pstricks}%
\usepackage{multido}%
\usepackage{pst-plot}%
\startdatakeys
    \source{24GENMATMEAG1}
    \BO{2026}
    \cycle{4}
    \competencesDuSocle{CA,RA}
    \objectifsApprentissage{C4A1aa,C4A2ad,C4B2cb,C4D2aa,C4D3ac}
    \motsCles{probabilité}
\enddatakeys
\begin{document}%
    


\medskip

Au casino, la roulette est un jeu de hasard pour lequel chaque joueur mise au choix sur un ou plusieurs numéros.

On lance une bille sur une roue qui tourne, numérotée de 0 à 36.

La bille a la même probabilité de s’arrêter sur chaque numéro.

\begin{center}
\psset{unit=1cm}
\def\ray{3cm}
\def\xmin {-4}  \def\xmax{4}
\def\ymin{-4}  \def\ymax{4}
\begin{pspicture}(\xmin,\ymin)(\xmax,\ymax)
%\psgrid[subgriddiv=1,gridlabels=0, gridcolor=lightgray] 
%\psaxes[arrowsize=3pt 3, ticksize=-2pt 2pt, labels=none](0,0)(\xmin,\ymin)(\xmax,\ymax)
\pscircle(0,0){\ray}
\multido{\na=90.00+19.46,\nb=99.73+19.46}{18}{
\pswedge[fillcolor=black,fillstyle=solid](0,0){\ray}{\na}{\nb}}
\psline(0,0)(\ray;80.27)
\pscircle[fillcolor=white,fillstyle=solid](0,0){1.5}
\pscircle(0,0){2.25} \pscircle[linewidth=5pt](0,0){3.3}
\bfseries
{\white
\uput[94.875]{4.875}(2.3;94.875){\scalebox{0.7}[1.5]{\textsf{26}}}
\uput[114.375]{24.375}(2.3;114.375){\scalebox{0.7}[1.5]{\textsf{35}}}
\uput[133.875]{43.875}(2.3;133.875){\scalebox{0.7}[1.5]{\textsf{28}}}
\uput[153.375]{63.375}(2.3;153.375){\scalebox{0.7}[1.5]{\textsf{29}}}
\uput[172.875]{82.875}(2.3;172.875){\scalebox{0.7}[1.5]{\textsf{22}}}
\uput[192.375]{102.375}(2.3;192.375){\scalebox{0.7}[1.5]{\textsf{31}}}
\uput[211.875]{121.875}(2.3;211.875){\scalebox{0.7}[1.5]{\textsf{20}}}
\uput[231.375]{141.375}(2.3;231.375){\scalebox{0.7}[1.5]{\textsf{33}}}
\uput[250.875]{160.875}(2.3;250.875){\scalebox{0.7}[1.5]{\textsf{24}}}
\uput[270.375]{180.375}(2.3;270.375){\scalebox{0.7}[1.5]{\textsf{10}}}
\uput[289.875]{199.875}(2.3;289.875){\scalebox{0.7}[1.5]{\textsf{8}}}
\uput[309.375]{219.375}(2.3;309.375){\scalebox{0.7}[1.5]{\textsf{11}}}
\uput[328.875]{238.875}(2.3;328.875){\scalebox{0.7}[1.5]{\textsf{13}}}
\uput[348.375]{258.375}(2.3;348.375){\scalebox{0.7}[1.5]{\textsf{6}}}
\uput[367.875]{277.875}(2.3;367.875){\scalebox{0.7}[1.5]{\textsf{17}}}
\uput[387.375]{297.375}(2.3;387.375){\scalebox{0.7}[1.5]{\textsf{2}}}
\uput[406.875]{316.875}(2.3;406.875){\scalebox{0.7}[1.5]{\textsf{4}}}
\uput[426.375]{336.375}(2.3;426.375){\scalebox{0.7}[1.5]{\textsf{15}}}
}
\uput[104.625]{14.625}(2.3;104.625){\scalebox{0.7}[1.5]{\textsf{3}}}
\uput[124.125]{34.125}(2.3;124.125){\scalebox{0.7}[1.5]{\textsf{12}}}
\uput[143.625]{53.625}(2.3;143.625){\scalebox{0.7}[1.5]{\textsf{7}}}
\uput[163.125]{73.125}(2.3;163.125){\scalebox{0.7}[1.5]{\textsf{18}}}
\uput[182.625]{92.625}(2.3;183.625){\scalebox{0.7}[1.5]{\textsf{9}}}
\uput[202.125]{112.125}(2.3;202.125){\scalebox{0.7}[1.5]{\textsf{14}}}
\uput[221.625]{131.625}(2.3;221.625){\scalebox{0.7}[1.5]{\textsf{1}}}
\uput[241.125]{151.125}(2.3;241.125){\scalebox{0.7}[1.5]{\textsf{16}}}
\uput[260.625]{170.625}(2.3;260.625){\scalebox{0.7}[1.5]{\textsf{5}}}
\uput[280.125]{190.125}(2.3;280.125){\scalebox{0.7}[1.5]{\textsf{23}}}
\uput[299.625]{209.625}(2.3;299.625){\scalebox{0.7}[1.5]{\textsf{30}}}
\uput[319.125]{229.125}(2.3;319.125){\scalebox{0.7}[1.5]{\textsf{36}}}
\uput[338.625]{248.625}(2.3;338.625){\scalebox{0.7}[1.5]{\textsf{27}}}
\uput[358.125]{268.125}(2.3;358.125){\scalebox{0.7}[1.5]{\textsf{34}}}
\uput[377.625]{287.625}(2.3;377.625){\scalebox{0.7}[1.5]{\textsf{25}}}
\uput[397.125]{307.125}(2.3;397.125){\scalebox{0.7}[1.5]{\textsf{21}}}
\uput[416.625]{326.625}(2.3;416.625){\scalebox{0.7}[1.5]{\textsf{19}}}
\uput[436.125]{346.125}(2.3;436.125){\scalebox{0.7}[1.5]{\textsf{32}}}
\uput[445.875]{355.875}(2.3;445.875){\scalebox{0.7}[1.5]{\textsf{0}}}
\psdots[dotscale=1.5](1.9;85.5)
\end{pspicture}
\end{center}

\medskip

\begin{enumerate}
\item Expliquer pourquoi la probabilité que la bille s’arrête sur le numéro 7 est $\dfrac{1}{37}$.
\item Déterminer la probabilité que la bille s’arrête sur une case à la fois noire et paire.
\item 
	\begin{enumerate}
		\item Déterminer la probabilité que la bille s’arrête sur un numéro inférieur ou égal à 6.
		\item En déduire la probabilité que la bille s’arrête sur un numéro supérieur ou égal à 7.
		\item Un joueur affirme qu’on a plus de 3 chances sur 4 d’obtenir un numéro supérieur ou égal à 7. A-t-il raison ?
	\end{enumerate}
\end{enumerate}

\bigskip
\showthesource
\end{document}
